\section{Introduction}

\begin{comment}
\fixme{Following could be possible format of the Introduction \begin{enumerate}
    \item A paragraph about the bit background about  = social media has been used by billions for communication + news. However, the problem is misinformation on social media. Some citations 

    \item What has been done in the past - Like a small paragraph. People have used blabla techniques to control misinformation. Any work using KG 

    \item What you have done in this thesis - We used KG technique + RL approach. Write a few lines about approach. Write 1 or 2 lines about dataset. Results 2 or 3 lines.

    \item Major highlights of this work are as following, then write the bullet points. You must highlight how this work is different from -previous works .. 
    
    \item Rest of the thesis is organized as follows.
    \begin{itemize}
        \item Chapter \ref{chp:Prelim} Preliminaries, discuss various terms which are useful for the readers?
        \item Chapter \ref{} ?? presents different related works related to the topic of the thesis.
        \item ??
    \end{itemize}
    
\end{enumerate}}
\end{comment}

Social media platforms are not only used by billions for communication but has become an alternative platform (compared to TV, print media) to receive news \cite{Boczkowski2017IncidentalNH}. These platforms have also being misused for spreading misinformation \cite{Flintham2018,jagtap2021misinformation} and its various forms such as fake news \cite{mayank2022deap,dhawan2022game}, rumors \cite{butt2022goes,sharma2021identifying}, misleading news \cite{sharma2022mis}, etc. %\Shakshi{Jagtap, Sabur, DEAP-FAKED, IJCNN, YouTube, GAME-ON, Misleading papers added} %However, this high level of communication has also brought forth a high number of misinformation spread on online social media \cite{Flintham2018}.
%As a result,
%To combat this various methods have been proposed to tackle this menace. 
With the sheer pace of misinformation generated manual fact-checking can not keep up. So, automated fact-checking %\st{provides a shelter} 
must be employed to tackle the menace of misinformation. One of the most prominent ways is through the usage of machine learning (ML) techniques \cite{islam2020deep,shakshi2023mis}. %\Shakshi{CIKM paper added} 
The majority of such techniques handle the context of the misinformation for its detection. For example, investigating what kind of users propagate the misinformation \cite{ghenai2018fake} or processing the content with natural language processing (NLP) techniques to derive clues from the way the misinformation is written, etc. \cite{su2020motivations}. % for its detection\par %Only a small proportion of misinformation detection works are concerned with checking the facts \textcolor{red}{@Gustav, kindly add a citation here if possible?}.



%\st{However, with such a} \fixme{In addition,} 
It is important to mention that the rise in misinformation has bring forth the focus on fact-checking as it helps in controlling the spread of misinformation \cite{martin2022facter,tymoshenko2021strong,sharma2022facov}. %\Shakshi{FaCov paper added}
 Providing explanations for fact-checking tools leads to trust among the people and offer better usability. Examples of the works related to explanations include formulating the task of explanation generation as an extractive-abstractive summarization task using transformer models \cite{Atanasova2020, Kotonya2020}. These particular works deal with specific domains such as political news \cite{Atanasova2020} and public health \cite{Kotonya2020}. In addition, to use these approaches, the datasets used in these studies have been annotated with gold-standard justifications by journalists. This approach is not scalable as annotation is a costly task. Alternatively, Knowledge graphs (KGs) have also been employed to provide a mechanism in which reliable third-party sources such as scientific articles and Wikipedia are used for fact-checking \cite{Ahmadi2019, GadElrab2019, Schatz2021}.  However, the use of reinforcement learning (RL) in such cases hasn't garnered much attention.%For example, in \cite{Sarrouti2021} health-related claims were evaluated against scientific articles using transformer models.


%to evaluate health-related claims against scientific articles using transformer models. KG 

%Sarrouti et al. \cite{Sarrouti2021} evaluated health-related claims against scientific articles using transformer models. Similar types  of works are done with KGs and explain the results of fact-checking \cite{Ahmadi2019, GadElrab2019, Schatz2021}.

%However, \fixme{can we add?- these solutions often lacks in} providing explanations for fact-checking methods. \fixme{Adding such a mechanism can} leads to trust and cooperation with the users and \st{leads to} \fixme{offer} better usability. Researchers \cite{Atanasova2020, Kotonya2020} have formulated the task of explanation generation as an extractive-abstractive summarization task using transformer models %introduced in \cite{liu2019}. 
%Atanasova et al. \cite{Atanasova2020} deals with political news and Kotonya et al. \cite{Kotonya2020} explores facts that require specific expertise with a public health-related dataset. %Both datasets are annotated with gold standard justifications by journalist. 
%The dataset used for training these methods approaches relies on the annotations which contain ground truth, unlike our method \fixme{we have not describe our method yet so let us not compared prevoius works with ours at this stage}, which extracts the truth from large-scale knowledge graphs.
%Sarrouti et al. \cite{Sarrouti2021} evaluated health-related claims against scientific articles using transformer models. Similar types  of works are done with KGs and explain the results of fact-checking \cite{Ahmadi2019, GadElrab2019, Schatz2021}. However, the use of reinforcement learning (RL) in such cases hasn't garnered much attention.\par

In this work, we propose an RL-based KG reasoning approach for explainable fact-checking that employs an RL agent to produce explanations in the form of relevant knowledge from the KG in the form of a path. The paths are arrived at by a technique from the field of KG reasoning known as multi-hop reasoning where an RL agent hops over facts in an attempt to connect two nodes %\fixme{can we write entities in brackets?} 
that correspond to real-world entities. The path describes the relation between those two entities %\fixme{entities or nodes?} 
and this path is used to classify a statement as true or false and also provides a human-readable explanation of the classification result. %We evaluated our approach using two widely used benchmark datasets, that is, FB15K-277 and NELL-995. %\fixme{Can we write a line or 2 for results?} 

To summarize, the main highlights of this work are as follows:
\begin{enumerate}
    \item To the best of our knowledge, this is the first work that has used a multi-hop reasoning agent for explainable fact-checking. We use two KGs (FB15K-277 and NELL-995) that are commonly used in multi-hop reasoning tasks to evaluate the performance of the proposed model. The KGs represent real-world relationships, sets of true facts are gathered from these KGs, and corresponding subsets of false facts are generated to simulate false facts.
    \item The results show that a multi-hop reasoning model can successfully check facts and provide highly human-readable explanations for the classifications. The agent is also able to provide explanations and a basis for classification which is not directly tied to any specific fact-checking task.
\end{enumerate}

The rest of the paper is organized as follows. Section \ref{related} discusses the related works. Next,  in Section \ref{methodology} we describe the methodology of the proposed approach, followed by the experiments and its results in Section \ref{experiments}. We conclude in Section \ref{conclusion} with some future directions. 